\chapter{Unravelling babel}

\begin{itemize}
    \item ``Now the whole earth had one language and the same words.'' So begins the story of the tower of Babel in the book of Genesis. 
    \item Most primate live in a pre-Babel world. A `leopard' call means the same thing to every vervet and green monkey from Senegal to South Africa \citep{Price2014}.
    \item This shows that animals don't have a language. Instead they have a system of signs (semiotics) where there is a one-to-one relationship between a given sign and a given meaning.
    \item That's how we think about language, but that's not how language works. \citet{Smith2019a} explain that when we communicate, whether using words or other methods like gestures, it's not just about the message itself. Instead, effective communication involves two key steps: showing and guessing. First, the person speaking or signalling tries to make their intentions clear (showing). Then, the person listening or observing uses these clues to understand what's being communicated (guessing). In this way, communication is like a game where one person has to think about the best way to get their message across, considering how the audience will interpret their words or actions, and the audience has to has to think about the best way to interpret them.
    \item In this view of language, humans never could have lived in a pre-Babel world because what we want to communicate is contingent. It depends on the communicative situation.
    \item That means that which constructions are grammatical depends on the communicative situation.
    \item The most general communicative situation, the kind of thing you assume when you write a book like this which could be read by anybody even well into the future is something like "English" or "French" or "Japanese". 
    \item What's grammatical in one may not be in another. This leads to dialectal differences and register differences.
    \item if two people share the context of "English" and "Japanese", they can mix constructions from both in the same sentence and it will be grammatical.
    \item A somewhat extreme, but not uncommon, case was when at the age of three our younger daughter, Inya, seemed to consider English a linguistic quirk of her father, restricted to com-sits with him alone. Her experience had been that Daddy spoke that way, and most people also used this way of speaking when interacting with him, but not in other situations. So, when our British in-laws came visiting, Inya refused to speak English to them because she found it awkward to use the ``Daddy register'' with anyone else -- even when it became clear they did not understand German. While this was an idiosyncratic identification of relevant com-sit characteristics, it in fact underlines the conventional and hence community-oriented nature of com-sits: such deviations put a spotlight on learner hypotheses that will be revised in later development, with more exposure to language use reflecting the conventional patterns, similarly to grammatical hypotheses from earlier acquisitional stages. \hfill\citep[42--43]{Wiese2023}
    \item Swearing feels right in some context and wrong in others. Perhaps this is because it belongs to some communicative situations -- it's grammatical there -- and not to others, where it's ungrammatical.
\end{itemize}






The ultimate point of the chapter, eventually, will be that a language is simply a communicative situation. When we're guessing how best to convey our intentions, the language is part of that. But languages are not just English, Japanese, Urdu, Apache, etc. They're much smaller -- swearing has different effects in different communicative situations/languages -- and they're much more mixed -- a bilingual conversation can mix languages within clauses and it will be fully grammatical. It all depends on what the sender knows and what they believe the audience knows.

It was 2003. We -- me, Yoko, and our toddler Tomo -- had just moved to Toronto from Tokyo. Bad timing, I know. That was the year of the first SARS outbreak, Toronto was SARS central, and  the city felt empty. Finding a teaching job was impossible when all the international students were understandably panicking.

Because of this, for several months, I was mostly at home with Yoko, our son Tomo, and my mom. We were lucky we could stay with her while we got our feet under us.

Yoko used Japanese exclusively with Tomo. I did too, except for when I'd read to him, which was always in English. My mom spoke only English, so he had a great bilingual environment and Tomo was picking up both Japanese and English as well as any toddler could.

I can't pinpoint when exactly, maybe six months after we arrived, but there was one day in the kitchen. Yoko and I were talking in Japanese, and Tomo and my mom were also there. During a pause, Tomo turned to her and started explaining what Yoko and I had been saying. I was amazed. It wasn't just that he knew there was a difference between Japanese and English; he also knew who understood what language.

\citet{Smith2019a} explain that when we communicate, whether using words or other methods like gestures, it's not just about the message itself. Instead, effective communication involves two key steps: showing and guessing. First, the person speaking or signalling tries to make their intentions clear (showing). Then, the person listening or observing uses these clues to understand what's being communicated (guessing). In this way, communication is like a game where one person has to think about the best way to get their message across, considering how the audience will interpret their words or actions, and the audience has to has to think about the best way to interpret them.

\subsection{More languages than you think}

Heike \citet{Wiese2023} tells a similar story about her daughter and communicative situations, what she calls \textsc{com-sits}.

\begin{quote}
    A somewhat extreme, but not uncommon, case was when at the age of three our younger daughter, Inya, seemed to consider English a linguistic quirk of her father, restricted to com-sits with him alone. Her experience had been that Daddy spoke that way, and most people also used this way of speaking when interacting with him, but not in other situations. So, when our British in-laws came visiting, Inya refused to speak English to them because she found it awkward to use the ``Daddy register'' with anyone else -- even when it became clear they did not understand German. While this was an idiosyncratic identification of relevant com-sit characteristics, it in fact underlines the conventional and hence community-oriented nature of com-sits: such deviations put a spotlight on learner hypotheses that will be revised in later development, with more exposure to language use reflecting the conventional patterns, similarly to grammatical hypotheses from earlier acquisitional stages. \hfill\citep[42--43]{Wiese2023}
\end{quote}

Is it a step too far to say that Inya experiences speaking to her grandparents as ungrammatical? 

\citet{Grieve2023} claims

\begin{quote}
    that languages can vary in their situational diversity. Languages that have been adapted for a wider range of communicative contexts are more situationally complex than languages that have been adapted for a narrower range ofcommunicative contexts. To support this claim, I consider examples of situational diversity from across a range of different languages and varieties of languages, drawing on empirical research from linguistics and anthropology. Second, I claim that situational diversity can help explain variation in grammatical complexity. I propose that increasing situational diversity in a language over time should lead to decreasing grammatical complexity. Furthermore, I argue that this trade-off between situational and grammatical complexity could explain how overall linguistic complexity could be maintained across languages and over time.
\end{quote}